\section{The worst exchange at a prescribed width}
\label{sec:exchange}
Let $N=\{1,\ldots,p\}$ index the input features. Fit an additive predictor
$b+\sum_{i\in N}a_i(x_i)$ and center each fitted component on the training
sample, absorbing its mean into $b$. On $n$ voting rows, write $a_i$ for
the matrix of contributions of feature $i$. Its columns correspond to
output coordinates. The target $y$ can be a scalar response, one-hot class
vector or full-model output. Set $r=y-b$ and use the inner product
$\langle a,c\rangle=n^{-1}\sum_{s=1}^n a_s^\top c_s$. For a support $S$,
define $U(S)=-\|r-\sum_{i\in S}a_i\|^2$ with respect to this inner product.
Expanding the squared norm gives
\begin{equation}
 U(S)-U(\varnothing)=\sum_{i\in S}d_i-2\sum_{\substack{i<j\\i,j\in S}}G_{ij},
 \qquad G_{ij}=\langle a_i,a_j\rangle,\qquad
 d_i=2\langle r,a_i\rangle-G_{ii}.
 \label{eq:game}
\end{equation}
The quadratic term arises from the shared prediction error even when the
predictor itself has only univariate terms. Its sign depends on how the fitted
components vary together on the voting rows.

Fix $1\le k<p$ and features $i\ne j$. For a background $T\subseteq
N\setminus\{i,j\}$ with $|T|=k-1$, write $T+i=T\cup\{i\}$ and
$\Delta_{ij}(T)=U(T+i)-U(T+j)$. At an improvement threshold $t$ this
background casts the vote $\mathbf 1\{\Delta_{ij}(T)>t\}$. Every background
votes for $i$ precisely when $t<D_{ij}(k)$, where
$D_{ij}(k)=\min_{|T|=k-1}\Delta_{ij}(T)$ with $i,j$ excluded. We use the
zero threshold. A positive minimum gives a strict preference at the target
width. Equation~\eqref{eq:game} lets us compute this decision without
sampling coalitions.

\begin{proposition}[Least favorable exchange]
For each $\ell\in N\setminus\{i,j\}$, set $z_\ell=G_{i\ell}-G_{j\ell}$.
Write $z_{(1)}\ge\cdots\ge z_{(p-2)}$ for the sorted differences. Then
\begin{equation}
 D_{ij}(k)=d_i-d_j-2\sum_{h=1}^{k-1}z_{(h)}.
 \label{eq:exchange}
\end{equation}
Any $k-1$ features achieving this upper sum form a minimizing background.
\end{proposition}
\begin{proof}
Subtract~\eqref{eq:game} for $T+j$ from the same expression for $T+i$.
All terms involving only $T$ cancel. The remaining gain is
$d_i-d_j-2\sum_{\ell\in T}z_\ell$. At fixed size the largest sum is attained
by the $k-1$ greatest differences, including negative differences when
necessary.
\end{proof}

Without a width constraint the worst background contains all positive
differences, giving $d_i-d_j-2\sum_\ell\max(z_\ell,0)$. Up to the choice
of minimization or maximization, this is the unrestricted comparison of
Ohno and Togawa~\cite[Eq.~19]{ohno2025posets}. For example, suppose
$d_i-d_j=7$ and the two remaining differences are $3$ and $2$. The
unrestricted lower bound is $-3$. At width two, the background can contain
only one more feature, and the minimum gain is $1$. Testing the unrestricted
bound first can settle a comparison without sorting. A full pass over $q$
currently free columns costs $O(q^3\log q)$ with pairwise sorting and needs
$O(q^2)$ storage once the Gram matrix is available.

\section{Feature fixing and the remaining search}
Write $i\succ_k j$ if $D_{ij}(k)>0$. Tied exchanges impose no constraint.
Any maximizer $S$ of $U$ at width $k$ satisfies $j\in S\Rightarrow i\in S$
for every such pair. Removing $j$ from a support lacking $i$ leaves a
background on which replacement strictly improves $U$.

\begin{proposition}[Positions in optimal supports]
The relation $\succ_k$ is a strict partial order. If at least $k$ features
dominate $j$, every maximizing support omits $j$. If $i$ dominates at least
$p-k$ features, every maximizing support contains $i$.
\end{proposition}
\begin{proof}
For transitivity assume $i\succ_k j\succ_k h$ and take a size-$(k-1)$
background $T$ excluding $i,h$. If $j\notin T$, replace $h$ first by $j$
and then by $i$ in $T+h$. If $j\in T$, first replace $j$ by $i$, then $h$
by $j$. Both routes strictly increase the utility and end at $T+i$.
Strict improvement rules out a cycle. A selected feature has space for at
most $k-1$ dominators beside it. When a feature is absent, at most
$p-k-1$ other features can also be absent, and these must include every
feature it dominates.
\end{proof}

These counts exhaust what the order alone implies about an individual
column. With dominators placed first, a feature $i$ can occur as early as
position $1+\#\{h:h\succ_k i\}$ in an ordering consistent with all the
comparisons. Placing the features dominated by $i$ last lets $i$ occur as
late as $p-\#\{j:i\succ_k j\}$. Taking the first $k$ entries of each
ordering gives the two extrema for inclusion. The utility values distinguish
the supports still allowed by these constraints.

After fixing a set $I$ of required features, each remaining feature $\ell$
gets the linear coefficient $d_\ell-2\sum_{h\in I}G_{\ell h}$ and the
remaining budget is $k-|I|$. The interaction coefficients within the free set
stay the same. We reapply~\eqref{eq:exchange} until no feature is fixed or
the residual width determines every choice. At this point enumeration can
skip any support that takes a feature without all its dominators. A binary
optimizer can impose the same inequalities $x_j\le x_i$ directly.
When the initial search varies only a window around a preliminary ranking,
these results concern the optimum conditional on the features fixed outside
that window.

This procedure fixes the component functions before optimizing the support.
Screening for $\ell_0$-regression with coefficients reestimated on the selected
variables solves a different objective~\cite{atamturk2020screening}. Once a
support has been chosen here, its additive predictor requires only $b$ and
the retained univariate functions. The number of inputs required at prediction
time is $k$. Fitting the original predictor may use all $p$ inputs.
